% FILE: 02_lineage_and_context.tex
% PROJECT: experiments-core
% THESIS ROLE: benchmark-experiments-and-equation-verification

\section{Lineage and Context}
\label{sec:experiments-core-lineage}

\subsection{2016 Language-Model Lesson}
\label{sec:experiments-core-lineage-2016}

The experiments in this directory originate from a 2016 observation that local text
imitation does not conserve consequence. A language model trained to reproduce surface
patterns of a corpus cannot be held to the causal structure of the process that
produced that corpus. The model may emit text that resembles a conforming trace without
having replayed any token through any net. This lesson established a foundational
requirement: any system that claims to reason about process must produce \emph{falsifiable
evidence of replay}, not plausible text. The Python validation scripts in this project
are the direct descendants of that lesson---they verify equations over concrete token
counts, not over surface-plausible narrative descriptions of conformance. The
\texttt{validate\_log\_fitness.py} script, for instance, computes fitness from raw
\texttt{(missing, consumed, remaining, produced)} tuples and asserts exact floating-point
results. No language model interpolation is involved at any step.

\subsection{Enterprise Process Gap}
\label{sec:experiments-core-lineage-enterprise}

The enterprise process gap---the observation that declared process models diverge
systematically from the processes that enterprises actually execute---motivates the
benchmark comparison between PM4Py and wasm4pm. PM4Py, as a Python library for
offline process mining, represents the incumbent tool for discovering and measuring
this gap. The pm4py-comparison experiments documented in
\texttt{pm4py-comparison/matrix.md} and \texttt{pm4py-comparison/zero-knowledge-benchmarks.md}
quantify the cost of remaining within that incumbent framework: DFG Discovery at
37.5$\times$ overhead, Alpha Miner at 34.2$\times$, Token Replay at 39.3$\times$,
and ZKP conformance timing out at $2^{30}+$ cycles versus wasm4pm's $2^{18}$ cycles.
These numbers are presented as empirical findings from the benchmark report; however,
as noted in the open questions (Section~\ref{sec:experiments-core-open}), no
reproducible benchmark harness with verifiable BPI Challenge 2017 log subsets is
currently present in the repository. The gap between declared and actual process is
therefore also present in the benchmark layer itself---a finding that the experiment
layer must eventually close.

\subsection{Relationship to the Chatman Equation}
\label{sec:experiments-core-lineage-chatman}

The Chatman Equation $A = \mu(O^*)$ asserts that alignment is a function of the
object-centric process state. The experiments in this directory operationalize that
equation at the measurement layer. The fitness equation
$f(L,N) = 1 - \frac{\sum \mathit{freq} \cdot m}{\sum \mathit{freq} \cdot c} - \frac{\sum \mathit{freq} \cdot r}{\sum \mathit{freq} \cdot p}$
is the Van der Aalst formalization of alignment under token replay semantics---it
measures precisely the degree to which an observed log $L$ conforms to a net $N$.
When this fitness is typed as \texttt{Metric<FITNESS, NUM, DEN>} with
\texttt{Between01<NUM, DEN>} as a compile-time type law, the Chatman Equation
becomes machine-checkable: a system whose compile-time state includes a valid
\texttt{Receipt} is a system that has, by construction, computed alignment within
the $[0,1]$ interval. The experiments verify both the mathematical and executable
sides of this correspondence. [AUTHOR THESIS: the full Chatman Equation unification
across all six process-intelligence products is not within the scope of the experiments
directory alone and is addressed in the dissertation synthesis chapter.]

\subsection{Relationship to Receipts and Replay}
\label{sec:experiments-core-lineage-receipts}

The receipt-and-replay architecture appears in this project through three artifacts.
First, \texttt{replay\_receipt\_sample.md} provides a \texttt{CryptographicReplayReceipt}
JSON schema with a concrete instance recording fitness~=~0.982 with BLAKE3 model
and log hashes and an Ed25519 signature. Second, \texttt{EVIDENCE\_CHAIN\_TRACE.md}
documents the 7-step chain: Raw OCEL $\to$ Parse $\to$ Admit $\to$ Token Replay /
Alignment $\to$ Metric $\to$ Receipt $\to$ Board Claim. Third,
\texttt{visualizer/blockchain.js} implements a SHA-256 cryptographic audit chain
with a genesis block and back-linked hash per event, demonstrating that the receipt
design is not confined to Rust---it is an architectural pattern that can be instantiated
in any runtime environment that supports a secure hash function. The OTel-BLAKE3
integration sample further shows that OpenTelemetry span data can feed directly into
the receipt chain using the formula
$R_i = \mathrm{BLAKE3}(R_{i-1} \;\|\; \mathrm{TraceID} \;\|\; \mathrm{SpanID} \;\|\; \mathrm{ParentSpanID} \;\|\; \mathrm{SpanName} \;\|\; \mathrm{StartTime} \;\|\; \mathrm{EndTime} \;\|\; \mathrm{InstructionCount})$,
making the audit chain extend from process-mining conformance all the way to hardware
instruction counts.
